% ----------------------------------------------------------
\chapter{Considerações Finais}\label{cap:consideracoesfinais}
% ----------------------------------------------------------

O presente trabalho teve como objetivo desenvolver um modelo classificatório de \textit{Machine Learning} para a identificação de indivíduos com perfil associado à hiperglicemia crônica clinicamente significativa, a partir de variáveis não invasivas da Subamostra Laboratorial da PNS 2013, visando um protótipo conceitual de triagem interpretável para a APS. Diante do problema de pesquisa, os resultados indicam que essa aplicação viabiliza um modelo classificatório e interpretável, porém com alcance delimitado. Os modelos alcançaram AUC-ROC entre 0,770 e 0,786 no conjunto de teste, e o modelo de referência, o \textit{Random Forest} com o conjunto de atributos compostos, obteve sensibilidade de 0,807 no ponto de operação principal, encaminhando a exame 40,5\% dos indivíduos avaliados, contra 84,0\% da regra clínica utilizada como comparador. Entretanto, o modelo de referência não cumpriu o critério de sucesso estabelecido antes da execução, avaliado no regime de sensibilidade equiparada, no qual as estratégias se mostraram praticamente equivalentes, e as suas saídas não constituem estimativas de risco individual, mas escores de ordenação, válidos para priorizar o encaminhamento por meio do limiar calibrado.

Quanto aos objetivos específicos, o levantamento de trabalhos correlatos fundamentou o referencial teórico e permitiu situar o desempenho obtido na mesma ordem de grandeza de modelos de identificação de diabetes não diagnosticado reportados na América do Sul, sem que essa aproximação constituísse uma comparação direta. A seleção de variáveis, orientada pelos princípios definidos na Metodologia, resultou em três conjuntos de atributos, dos quais o conjunto completo foi descartado em virtude da redundância entre o escore de multimorbidade e os seus componentes, e a RCEst foi excluída dos demais em virtude da colinearidade com a circunferência da cintura. Na comparação entre os algoritmos, o \textit{Random Forest} superou a Regressão Logística em sensibilidade no ponto de operação principal, detectando um superconjunto dos casos detectados pela regressão, com diferença de acerto global não significativa, enquanto a sua PR-AUC inferior concentrou-se nos indivíduos de maior escore. Essa vantagem, contudo, não pôde ser atribuída à captura de relações não lineares entre as variáveis, visto que nenhum experimento foi desenhado para isolar esse efeito. Por fim, a avaliação incluiu, além das métricas previstas, a PR-AUC, o NNS, a taxa de encaminhamento e a calibração das probabilidades, e a explicabilidade por meio do SHAP foi implementada sobre o \textit{Random Forest}, permanecendo a explicabilidade da Regressão Logística como etapa não executada.

Quanto ao objetivo geral, o trabalho entrega a classificação individual explicada, isto é, a decisão de encaminhamento acompanhada da decomposição das contribuições de cada variável para o indivíduo, e não uma estimativa de probabilidade individual. Essa decomposição foi demonstrada nos casos do conjunto de teste e apresentada no painel de resultados, desenvolvido como artefato de apresentação. A API construída com o \textit{framework} FastAPI estende essa entrega a novos indivíduos, retornando, para os dados informados no formulário da página de classificação individual do painel, a decisão de triagem, o escore de ordenação em relação ao limiar e a decomposição das contribuições, acompanhada da restrição de leitura.

Como contribuições, destacam-se quatro aspectos. O primeiro é a avaliação operacional dos modelos pela taxa de encaminhamento e pelo NNS, com o regime de sensibilidade equiparada, que tornou explícita a troca entre casos detectados e carga gerada para a rede de saúde. O segundo é a avaliação da calibração das probabilidades, aspecto que a revisão consultada apontou como ausente nos estudos de triagem de diabetes realizados na América do Sul, e que delimitou a forma correta de leitura das saídas. O terceiro é a estratégia construída para medir a exposição das explicações à multicolinearidade, por meio do desenho dos conjuntos de atributos, da execução em dois regimes e do reporte dos pares e blocos correlacionados. O quarto é a reprodutibilidade do trabalho, assegurada por um comando único que gera todos os resultados a partir dos microdados, por conferências automáticas contra tabelas de referência, por testes automatizados e pelo registro das decisões de implementação, de maneira que o percurso entre o projeto de pesquisa e os resultados pode ser reconstituído decisão a decisão.

Como trabalhos futuros, sugere-se, em primeiro lugar, a comparação entre os modelos treinados com e sem reponderamento das classes, seguida da recalibração das probabilidades, caso o reponderamento se mostre vantajoso para a discriminação. Sugere-se também a avaliação em múltiplas partições ou com intervalos de confiança, a fim de dimensionar a incerteza das diferenças observadas, e a validação em uma base de dados externa. Além disso, permanecem como etapas a explicabilidade da Regressão Logística, o teste da hipótese de que a ausência da informação de colesterol esteja associada ao desfecho e a avaliação da estabilidade das categorias de raça/cor com poucas observações. Por fim, a integração do modelo a sistemas de informação da APS, como o PEC, depende da compatibilização entre a codificação das variáveis da PNS e a dos registros clínicos, etapa que antecede qualquer uso do modelo em serviço.
